\section{Summary of Contributions}
\label{sec:conclusion_summary_of_contributions}

This section revisits the four research questions introduced in Section~\ref{sec:introduction_research_questions}.
For each of them, it summarizes the answer provided by this thesis.

Research Question~\ref{rq:latent_space} concerned the approaches proposed in the literature to analyze the internal workings of \glspl{nn}.
To answer it, we first made explicit the working hypothesis shared by most of this literature, the \ConceptHypothesisName.
According to this hypothesis, \glspl{nn} generalize by organizing their \glspl{ls} around concepts.
We then organized existing methods into four families: probing, interpretable directions, \glspl{sae}, and clustering.

Research Question~\ref{rq:taxonomy} concerned the organization of existing \gls{xrl} methods \parencite{alaarabiou:hal-04685607, alaarabiou2026taxonomy}.
Our starting point was to clarify what an explanation means in the context of \gls{rl}.
Building on this framework, we identified the notions of subject and source of an explanation.
Regarding the subject, we distinguish two levels: a single decision and the agent's strategy.
For the source, we identify three possibilities: policies, trajectories, and value functions.
From these elements, we proposed a taxonomy that accounts for the specificities of \gls{rl} and used it to review the state of the art.

Research Question~\ref{rq:mechanisms} concerned the use of \glspl{lr} to identify the mechanisms underlying the decisions of \gls{rl} agents.
To address it, we introduced \gls{car}, an \gls{xrl} method based on the principles of \gls{me} \parencite{alaarabiou2026car}.
\gls{car} trains several surrogate policies independently to reproduce the same agent.
It then clusters their \glspl{lr} and compares the resulting partitions.
This procedure relies on the \RepresentationConvergenceHypothesisName, according to which networks trained on the same task develop similar concepts.
Applied to six Atari environments, \gls{car} recovers similar clustering structures across surrogate policies.
The qualitative analysis shows that these clusters capture interpretable aspects of the task.

Research Question~\ref{rq:evaluation} concerned the objective quantification of explanation quality.
We addressed it by introducing three metrics.
The \gls{cu} measures the consistency of the clusters across surrogate policies.
The \gls{as} measures how much these clusters differ from the structure of the observations, the actions, and an untrained network.
The \gls{acu} combines both properties into a single measure and provides a practical criterion for selecting the layer to analyze.

Throughout these contributions, we developed three hypotheses that form a coherent chain.
The \MechanisticHypothesisName states that \glspl{nn} can be understood through their internal mechanisms.
The \ConceptHypothesisName states that these mechanisms are organized around concepts.
The \RepresentationConvergenceHypothesisName states that these concepts reflect the task.
\gls{car} translates these hypotheses into an empirical procedure, and the experimental results provide support for this framework.
